\begin{abstract}
Las heladas meteorológicas son uno de los eventos climáticos extremos de mayor
impacto sobre la agricultura altoandina del Perú, generando pérdidas severas en
cultivos y ganado. Su predicción temprana constituye un problema de
clasificación binaria fuertemente desbalanceado, en el que los clasificadores
convencionales favorecen la clase mayoritaria. A diferencia de trabajos previos
que abordan este desbalance mediante sobremuestreo sintético (SMOTE) y suelen
omitir la corrección de anomalías de instrumentación, este trabajo propone un
pipeline de minería de datos que combina corrección explícita de errores de
sensor, ponderación balanceada de clases y optimización orientada a la clase
minoritaria. Sobre 1\,672\,827 observaciones meteorológicas horarias de
estaciones peruanas (18 variables), se identificó y corrigió un código de error
de sensor (\texttt{-999}) no documentado, se definió la variable objetivo
mediante el umbral físico de $0^{\circ}\text{C}$ y se eliminó el predictor
definitorio para evitar fuga de información. El preprocesamiento aplicó
imputación por mediana y \texttt{RobustScaler}, seguido de PCA que redujo el
espacio de 15 a 5 componentes reteniendo el 95\% de la varianza. Se entrenó un
\emph{Random Forest Classifier} balanceado, optimizado mediante
\texttt{GridSearchCV} (3-fold, F1-Score Macro). El modelo alcanzó un
\emph{Recall} de 92.34\% y un \emph{F1-Score} macro de validación de 0.8330 en
la detección de heladas, con un índice Kappa de Cohen de 0.6710, evidenciando
concordancia sustancial. Los resultados validan el sistema como herramienta de
alerta temprana de alta sensibilidad para la prevención de desastres
agrometeorológicos.

\vspace{0.5cm}
\noindent \textbf{Palabras clave:} 
Agriculture, data mining, meteorology, predictive models, principal component
analysis, random forests.
\end{abstract}